\selectlanguage{english}
\noindent\textbf{Chapter XII, section 5, manuscript pages 115--169.}
\medskip

These pages develop a concrete calculus for deciding when an ordinary gluing diagram of categories represents a homotopy gluing. Its central object is a comparison map. Given a strictly commuting square with initial corner $X_0$, adjacent corners $X_1,X_2$, and terminal corner $Y$, Grothendieck constructs a category $\mathcal X_Q$ from the three-category span and a canonical functor
\[
                    \varphi_Q:\mathcal X_Q\longrightarrow Y.
\]
The square is $W$-homotopically cocartesian when this particular functor belongs to the chosen weak-equivalence class $W$. The definition remembers the cocone, not only its four vertex homotopy types. The work that follows explains why the construction is independent of its two orientations, how smoothness can test the comparison on fibres, and why monomorphisms of presheaves provide a reliable supply of such squares.

The opening argument concerns a different but related exactness problem. Total $W$-asphericity here is the requirement that, for every $p:X\to Y$, the category of elements of the ordinary left Kan extension $p_!1_X$ be $W$-aspherical. It is the definition of this chapter, explicitly distinguished by the editors from the earlier use of the same phrase in \emph{Pursuing Stacks}. $W_0$-properness, or $W$-properness under W(5), allows Grothendieck to compute
\[
 p_!1_X\cong\operatorname*{colim}_{i\in I}p_{i!}1_{X_i}
\]
for a family $X\to I$. A terminal object of $I$ reduces this colimit to the last fibre. The commentary supplies the generators-and-relations proof on components and the precise functor on comma categories requested in footnote~83. It also proves that total $W_\infty$-asphericity is not invariant under weak homotopy equivalence. The properness hypothesis controls ordinary Kan extensions; equivalence of classifying spaces alone does not do that work.

The smooth integral of $X_1\leftarrow X_0\to X_2$ keeps the three sets of objects disjoint. Its mixed arrows run from the outer fibres to the middle fibre, with
\[
 \operatorname{Hom}_{\mathcal X_Q}(a,x)
       =\operatorname{Hom}_{X_i}(a,\alpha_i x)
       \qquad(a\in X_i,\ x\in X_0).
\]
Its universal property therefore includes natural transformations $f_i\alpha_i\Rightarrow f_0$. They need not be invertible. Strict commutativity of the given square supplies identity transformations and hence the canonical $\varphi_Q$. Dividing the integral into two full closed pieces gives a genuine categorical pushout and pullback. Each piece retracts onto its outer fibre by an adjunction. The nerve of this closed cover proves directly that the total category represents the homotopy pushout of classifying spaces. The general theorem is Thomason's; the bibliography distinguishes verification of the original metadata from verification of an independent primary proof \citep{Thomason1979,Cegarra2020}.

Even the smallest examples explain why this construction is needed. For the span $e\leftarrow\{a,b\}\to e$, the integral is a four-object category whose nerve is a circle. Its ordinary categorical pushout is a point. The ordinary pushout of the same diagram in the homotopy category is also a point: it retains equality of components but loses the two paths implementing compatibility. The integral retains those paths as transition arrows. This is the concrete content, on these pages, of the insistence on forming diagram categories before localization.

The excision hypotheses require particular care. Proposition~13 and its cube corollary are valid conditional assertions under the full W(7 bis) and W(7) printed in the manuscript. Those axioms include reflection statements stronger than the usual preservation of weak equivalences by homotopy pushout. They fail for $W_\infty$. The commentary proves this with a connected acyclic finite complex having nontrivial fundamental group. Its suspension is contractible, so the union of its two cones and both cones are contractible although their intersection is not. A map of this closed-cover square to the point square violates the converse cube implication. The usual forward gluing theorem survives. A normalized-chain calculation also explains exactly why the complete three-way implication works for integral homology equivalences. The earlier chapter's claim about full W(7) for $W_\infty$ is therefore not a valid justification of the converses in the present portion.

The proper integral reverses the transition arrows. Grothendieck compares the two constructions by the functorial resolution
\[
                         C\longleftarrow\Sigma(C)\longrightarrow C^{\mathrm{op}}.
\]
Both projections are aspherical, and the same comparison can be integrated. Proposition~14 proves agreement of the four versions obtained from smooth or proper integration and from the span or its opposite. The commentary writes out the chosen twisted-arrow convention and completes the deferred transitivity verification. The natural comparison with the opposite operation is involutive; its naturality makes the square of four integration comparisons commute. The indexing category must keep its two paths distinct until they reach $Y$. Replacing that free diamond by a commutative square would erase precisely the information the proof is using.

The question on manuscript pages 145--146 asks whether a bad ordinary pushout can become a good homotopy cocone after mapping to another target. We answer it with finite posets. A collapsible disk with seven triangles maps to the tetrahedral two-sphere so that all associated comma categories are nonempty and connected. Its face poset $E$ embeds as a sieve in a contractible mapping-cylinder poset $B$. The categorical pushout
\[
                         P=B\amalg_E e
\]
has two incomparable minimal vertices below the face poset of that sphere. Its nerve is their join with the subdivided sphere, hence has classifying space $S^3$. The span $B\leftarrow E\to e$ has contractible integral. Consequently the comparison to $P$ is not a weak equivalence, whereas its composite with $P\to e$ is. The proof includes the categorical pushout universal property; a separate exact finite calculation checks the collapses, comma connectivity, and rational homology. This is a worked development, with no claim of priority. The same disk-to-sphere map supplies the promised counterexample to homotopy invariance of total asphericity.

Theorem~4 turns from construction to recognition. Its forward assertion says that smoothness of the three maps $f_i$ implies smoothness of $\varphi_Q$. The reverse assertion for the middle map retains strong W(7 bis). Two slice calculations identify the strict fibres relevant to local smoothness with integrals of lift categories. Footnote~107 correctly observes that the apparent square in the original proof is not strictly commutative. The repair uses the split fibration of the integral over the fork and its aspherical fibres. This proves the forward direction without the strong excision axiom. An explicit contraction removes the same unnecessary hypothesis from the open-cover corollaries. The final unconditional sentence after Corollary~6, however, omits essential outer smoothness assumptions; a four-object lift category with an isolated component gives a finite counterexample.

The last part returns to ordinary category theory. A full sieve $U\subset X$ and its complement $T$ are glued by the presheaves $\operatorname{Hom}_X(-,t)$ on $U$. Proposition~15 says that a square along $U\subset X$ is a pushout exactly when it identifies the complementary categories and transports these crossing-arrow presheaves by left Kan extension. The connected factorization categories in its statement are the fibres of that comparison of presheaves. This gives a full proof of the criterion and explains why an agreement on objects and strata would be insufficient.

For a presheaf $F$ on $X$, the formula
\[
 N(X/F)_n=\coprod_{x_0\to\cdots\to x_n}F(x_n)
\]
then makes the final gluing theorem transparent. It proves preservation of colimits by the category-of-elements construction; the relative construction into $\Cat/X$ also preserves all small limits, the known strengthening recorded by \citet[Proposition~3.2.12 and Corollary~3.2.13]{Cisinski2006}. A pushout of presheaves along a monomorphism becomes a simplicial pushout along a monomorphism. Thus the positive homotopy-pushout clause of Theorem~5 holds for every fundamental localizer. Its subsequent reflection implications still require the strong hypothesis. Model structures require further conditions, supplied later by accessibility and local test categories; attachment stability alone is insufficient.

The later derivator definition identifies $\varphi_Q$ with the terminal component of the coherent left-Kan-extension counit \citep[Definition~3.9]{Groth2013}. The comparison with $\infty$-categories requires another localization: a total category models the categorical colimit after its transport arrows are inverted, while classifying spaces invert all arrows \citep[Corollary~3.3.4.3]{Lurie2009HTT}. These formulations clarify the target of the present calculations without replacing their elementary categorical proofs. The equivalence of the two integration conventions and the strict-cocone question are resolved, as are the specified coherence and local-smoothness gaps. The opening request for a sufficiently large class of cofibrations closed under the stated interval operation is not settled here in its full generality. We record that limitation without declaring it a present-day open problem.
